\begin{abstract}
本文提出原型路由图提示方法 \method，利用支持集原型将可学习提示分配给不同节点。在每类仅有少量标注的条件下，任务适配不仅需要学习提示内容，还需要确定提示应当作用于哪些节点。\method 从聚合节点属性构造类别参照，在保留支持集锚点的同时利用未标注图特征进行校正，并将节点与原型的关联转换为固定路由矩阵。该矩阵组合可学习类别提示，将其注入冻结图卷积编码器，最终采用原型分类。我们在 Cora、CiteSeer、PubMed 上分别评估 1、3、5-shot 设置，对各条件使用配对的支持抽样与训练种子。5-shot 下，三个数据集的 Accuracy 分别为 73.02\%、66.17\%、70.98\%；相对同配方共享提示，原型路由在对应条件下分别带来 36.94、40.45、3.02 个百分点的提升。这些逐条件配对结果支持原型关系对提示分配的组织价值，结论限于所评估的传导式节点分类设置。
\end{abstract}

\section{引言}
图提示通过少量任务相关参数适配预训练图神经网络，为少样本节点分类提供了一种可行路径。在这一任务中，每个类别仅有少量标注节点，需要尽可能利用已有表示完成任务适配。现有图提示方法可以在输入特征空间中进行调整、引入提示结构，或改变冻结表示的读出方式~\citep{gpf2023,allinone2023,graphprompt2023}。除提示注入位置之外，一个直接影响适配方式的问题是：提示的贡献应如何分配给不同节点？

共享提示在注入位置对所有节点施加相同的修改，能够表达全图一致的任务调整，却不依据节点与目标类别的关系分配提示。输入相关的提示组合提供了更灵活的选择~\citep{gpf2023}。与此同时，少量支持样本本身就给出了当前任务的类别参照。我们关注是否能够直接利用这些参照组织提示分配，使参数学习集中于提示内容，而不必额外学习一个路由网络。

本文的设计假设是，支持集原型可以作为节点适配的任务坐标。与同一类别关系相近的节点能够接收相关的提示组合，不同类别关系则对应不同的适配方向。这些原型无需成为模型的可学习参数，可以在提示优化之前由图属性和支持标签计算得到。由此，任务相关的提示分配与提示内容优化形成职责明确的两个部分。

基于这一思路，\method 首先对两跳聚合属性进行全图中心化和逐行归一化，以支持节点初始化类别原型，再利用未标注节点特征进行支持集锚定的校正。节点与原型的关联经过全图中心化和邻域聚合，形成固定路由矩阵。路由矩阵组合少量可学习类别提示，将其注入冻结 GCN 的第一隐藏层，最终采用原型分类。图~\ref{fig:architecture} 概括了固定路由与可学习提示的分工。

实验使用两组互补证据。主比较包括冻结预训练编码器、全量微调，以及在统一节点分类协议下实现的三类提示算子。路由消融保持提示参数化、注入位置、优化配方和读出不变，以共享组合替换节点相关路由；无提示条件用于衡量完整适配过程的作用。Cora、CiteSeer、PubMed 的各个 shot 条件分别报告，均值仅由同一设置下的不同 seed 计算。

本文的主要贡献如下：
\begin{itemize}
\item 提出原型路由提示构造，将支持集定义的类别关系转换为固定的节点相关系数，用于适配冻结图编码器。
\item 给出仅包含 $C\times128$ 个可学习提示参数的实现，无需可学习路由网络，并明确区分未标注特征使用与基于标签的拟合过程。
\item 在三个引文网络的 1、3、5-shot 条件下提供配对主比较和路由消融，逐条件报告 seed 波动及配对不确定性区间。
\end{itemize}
本文贡献聚焦类别参照对隐藏层提示的组织作用，不将原型、未标注原型校正或节点相关提示各自视为全新概念。

\section{相关工作}
\paragraph{图提示与任务适配。}
GPPT 研究图预训练与提示调优的下游适配范式~\citep{gppt2022}。GraphPrompt 以基于相似度的共同模板组织预训练和下游预测，并对任务读出进行提示适配~\citep{graphprompt2023}。All in One 从提示向量、提示结构及插入方式出发构建多任务图提示框架~\citep{allinone2023}。GPF 与 GPF-plus 在输入特征空间中引入提示，其中包括节点相关的提示组合~\citep{gpf2023}。本文区别在于利用支持集锚定的原型关系，预先计算隐藏层提示的固定分配方式；节点相关性本身不是本文的创新主张，也不据此声称对任意预训练目标均适用。

\paragraph{原型与未标注样本。}
Prototypical Networks 通过样本与支持集类别代表之间的距离进行少样本分类~\citep{proto2017}；半监督少样本学习也已利用未标注样本改善原型构造~\citep{semiproto2018}。因此，本文将原型理解为任务参照，重点研究其组织提示更新的作用。路由原型处于聚合属性空间，在提示训练中保持固定；最终预测原型则由适配后的编码器表示构造。本文评估同一图上的已知类别，不涉及通过元学习向未见类别迁移。

\paragraph{评估协议。}
ProG 为图提示方法提供统一评估框架与实现工具~\citep{prog2024}。本文采用与 ProG 兼容的 GCN 和固定边预测预训练权重，明确区分原论文方法与本项目统一协议下的实现。比较结果用于描述这一受控设置，不等同于复现各基线的最佳已发表配置。

\section{问题定义}
给定属性图 $G=(V,E,X)$，节点数为 $n$，特征矩阵为 $X\in\mathbb R^{n\times d}$，任务类别数为 $C$。每个类别 $c$ 的支持集 $S_c$ 包含 $k$ 个标注节点，其中 $k\in\{1,3,5\}$。记 $S=\bigcup_c S_c$，查询集 $T$ 与 $S$ 不相交。完整图结构与所有节点特征可见，拟合和超参数选择仅接收 $S$ 内的标签；查询标签在预测封存后用于评分。

编码器采用预训练两层 GCN~\citep{gcn2017}，隐藏维度与输出维度均为 128，在 \method 中保持冻结。本文属于同图、已知类别的传导式分类：查询特征参与图聚合、原型校正和全图中心化，不属于归纳式未见类别评估。文中的一个\emph{实验条件}固定数据集、shot 数、方法及预训练权重。

\section{原型路由图提示}
\label{sec:method}
\begin{figure}[t]
\centering
\fbox{\begin{minipage}[c][52mm][c]{0.92\linewidth}
\centering
\lang{Architecture figure placeholder}{主架构图占位符}\\[8mm]
\lang{Graph attributes and support labels $\longrightarrow$ anchored prototypes $\longrightarrow$ fixed routing}{图属性与支持标签 $\longrightarrow$ 支持集锚定原型 $\longrightarrow$ 固定路由}\\[4mm]
\lang{Fixed routing + trainable prompt tokens $\longrightarrow$ frozen GCN $\longrightarrow$ prototype prediction}{固定路由与可学习提示 $\longrightarrow$ 冻结 GCN $\longrightarrow$ 原型分类}
\end{minipage}}
\caption{\lang{Overview of \method. Support labels initialize class references in aggregated attribute space. Unlabeled graph features refine these references, and the resulting node--prototype associations determine a fixed routing matrix. Only the class prompt tokens are optimized; the routed prompts enter the first hidden layer of a frozen GCN. Query labels are excluded from fitting and selection. The final architecture artwork is intentionally left as a placeholder.}{\method 总体流程。支持标签在聚合属性空间中初始化类别参照，未标注图特征用于校正参照，节点与原型的关联决定固定路由矩阵。仅类别提示向量参与优化，路由后的提示注入冻结 GCN 的第一隐藏层。拟合与选参不使用查询标签。主架构图按要求暂用占位符。}}
\label{fig:architecture}
\end{figure}

\subsection{支持集锚定的类别参照}
令 $\widehat A=\widetilde D^{-1/2}(A+I)\widetilde D^{-1/2}$ 为带自环的归一化邻接算子。由两跳聚合属性构造路由表示：
\begin{equation}
R=\rownorm\!\left(\widehat A^2X-\mathbf1\,\mean_{v\in V}(\widehat A^2X)_v\right).
\label{eq:attributes}
\end{equation}
其中 $\rownorm$ 对每一行进行欧氏归一化。路由使用图属性及结构，不依赖额外学习的网络。类别初始参照为
\begin{equation}
\mu_c^{(0)}=\frac{1}{k}\sum_{i\in S_c}R_i.
\label{eq:anchor}
\end{equation}

为利用未标注图信息并保留标注参照，记 $U_V=V\setminus S$。从 $\mu^{(0)}$ 开始进行三轮固定校正：
\begin{align}
a_{ic}^{(t)}&=\frac{\exp(R_i^\top\norm{\mu_c^{(t)}}/0.05)}{\sum_{c'}\exp(R_i^\top\norm{\mu_{c'}^{(t)}}/0.05)},\quad i\in U_V,\label{eq:assignment}\\
\mu_c^{(t+1)}&=0.5\mu_c^{(0)}+0.5\frac{\sum_{i\in U_V}a_{ic}^{(t)}R_i}{\max(\sum_{i\in U_V}a_{ic}^{(t)},10^{-8})},\quad t=0,1,2.\label{eq:refinement}
\end{align}
每轮锚点始终为最初的支持原型，而不是上一轮原型。这是固定的原型构造过程，本文不赋予其期望最大化目标或收敛保证。计算不使用 $U_V$ 中节点的标签。

\subsection{从类别关系到固定路由}
以校正后的原型计算温度为 $\tau_r=0.1$ 的节点—类别关联：
\begin{equation}
Q_{ic}=\softmax_c\!\left(R_i^\top\norm{\mu_c^{(3)}}/\tau_r\right),\qquad
B=\widehat A\left(Q-\mathbf1\,\mean_{v\in V}Q_v\right).
\label{eq:routing}
\end{equation}
中心化表示节点类别关联相对全图平均关联的偏离，邻域聚合进一步整合相邻节点的关联模式。$Q$ 是按行归一化的关联矩阵，$B$ 则包含可为负的提示组合系数，不能解释为最终类别概率。每个支持集仅计算一次 $B$，提示训练期间保持不变。

\subsection{冻结隐藏计算中的提示适配}
令 $H^{\mathrm{pre}}$ 为原图经过 GCN 第一层得到的预激活表示，学习零初始化的矩阵 $P\in\mathbb R^{C\times128}$。记 $s$ 为 $H^{\mathrm{pre}}$ 的逐通道样本标准差，下界设为 $10^{-6}$，则注入后的隐藏表示为
\begin{equation}
\widetilde H^{\mathrm{pre}}=H^{\mathrm{pre}}+10\,[B\tanh(P)]\odot s.
\label{eq:injection}
\end{equation}
尺度向量按节点广播。经过原有 LeakyReLU 与第二层图卷积后得到 $Z$。支持监督优化提示内容，路由与编码器参数保持固定；零初始化使训练起点恢复冻结编码器的原始计算。Cora、CiteSeer、PubMed 分别包含 896、768、384 个可学习提示参数。

\subsection{训练与预测}
使用 Adam 优化提示矩阵 100 步，学习率为 0.01，权重衰减为 0：
\begin{equation}
\mathcal L(P)=\frac{1}{|S|}\sum_{i\in S}\operatorname{CE}(\ell_i(P),y_i)
+10^{-3}\mean(P\odot P).
\label{eq:loss}
\end{equation}
支持 logits $\ell_i$ 采用温度为 0.2 的余弦原型分类。3-shot 和 5-shot 训练时，构造某个支持节点自身类别的原型需排除该节点；1-shot 使用包含自身的原型，因为排除唯一支持样本会导致该类别原型无法定义。预测时，对全图输出减去全图均值并逐行归一化，再由支持原型分类。上述读出细节在路由消融中保持固定，复现定义见附录~\ref{app:protocol}。

路由只需预先计算一次。在稀疏图上，两次属性聚合的复杂度为 $O(|E|d)$，每轮原型校正为 $O(nCd)$，每步提示组合 $B\tanh(P)$ 为 $O(nC\cdot128)$，此外还包含冻结编码器的前向与反向计算。冻结权重减少可学习参数，但不会消除经过第二层卷积的梯度计算。因此本文报告参数量，不声称未经测量的运行时间优势。

\section{实验}
\subsection{协议与比较方法}
\label{sec:protocol}
在 Cora、CiteSeer、PubMed 上分别评估每类 1、3、5 个支持节点。官方训练掩码定义支持池，官方测试掩码定义固定查询集。同一次支持抽样使用相同的逐类排列，取前 1、3、5 个节点构成嵌套支持集。每个条件包含 20 个支持抽样 seed 与 3 个训练 seed；各方法共享图、预训练权重及支持／查询身份。每个数据集采用项目中原始的边预测预训练权重，其预训练 seed 未记录。

主比较包含：\textbf{冻结预训练}，不添加提示并采用与 \method 相同的最终原型读出；\textbf{全量微调}，更新预训练编码器及任务读出；\textbf{GPF} 与 \textbf{GPF-plus}，在输入特征空间中适配；\textbf{GraphPrompt$^{\dagger}$}，统一节点协议下的对角读出提示实现。GPF-plus 使用 20 个提示向量。上标表示协议适配，不代表完整复现 GraphPrompt 原始的子图处理流程。这些基线遵循相应论文的提示算子思路~\citep{gpf2023,graphprompt2023}，但使用本文的编码器与选参规则。

\method 与共享提示消融使用同一固定配方。需要选参的比较方法在 3-shot 和 5-shot 下使用三个支持内折，候选为学习率 $\{0.001,0.01\}$ 与权重衰减 $\{0,0.0001\}$ 的四种组合。各候选拟合及最终拟合均为 100 步；1-shot 使用固定参数。因此，各方法的总优化预算并不完全相同，其余设置见附录~\ref{app:protocol}。

\paragraph{统计单位。}
在固定条件内，先对同一支持抽样下的 3 个训练 seed 求均值，再对 20 个支持抽样结果计算均值与样本标准差。不跨数据集或 shot 汇总平均表现。配对区间以支持抽样为块进行 10,000 次 bootstrap，保留方法之间的配对关系；报告固定图和权重条件下、未经多重比较校正的 95\% 区间。正文报告 Accuracy 与 Macro-F1，其余指标见附录~\ref{app:more}。

\subsection{相对预训练、微调与其他提示的适配表现}
主比较检验完整原型路由适配在相同图、权重及支持标签预算下能否改善预测。表~\ref{tab:main-accuracy} 与表~\ref{tab:main-macro_f1} 展示全部条件，不包含跨条件平均行。
\begin{table}[!htbp]
\centering
\caption{\lang{Accuracy (\%, higher is better). Main comparison. Within each dataset/shot/method, training seeds are averaged within each support draw; mean and sample standard deviation are then computed over 20 support draws. Bold marks the largest mean in that row. $\dagger$: common-protocol adaptation.}{Accuracy（\%，越高越好）。主比较。固定数据集、shot 与方法，在每个支持抽样内先平均训练 seed，再对 20 个支持抽样计算均值与样本标准差。粗体为该行最大均值。 $\dagger$：统一协议适配实现。}}
\label{tab:main-accuracy}
\small
\setlength{\tabcolsep}{4pt}
\resizebox{\linewidth}{!}{%
\begin{tabular}{lrrrrrr}
\toprule
\lang{Dataset / shot}{数据集 / shot} & \frozenname & \ftname & GPF & GPF-plus & \gpromptname & \method \\
\midrule
Cora / 1 & \meanstd{23.17}{5.48} & \meanstd{45.41}{10.18} & \meanstd{15.20}{3.24} & \meanstd{30.80}{6.64} & \meanstd{21.88}{4.62} & \bestmeanstd{54.59}{9.18} \\
Cora / 3 & \meanstd{29.42}{3.47} & \meanstd{61.54}{6.78} & \meanstd{24.32}{4.66} & \meanstd{40.21}{4.80} & \meanstd{27.75}{3.64} & \bestmeanstd{70.76}{2.76} \\
Cora / 5 & \meanstd{33.56}{5.03} & \meanstd{67.50}{5.06} & \meanstd{28.17}{5.54} & \meanstd{53.14}{5.55} & \meanstd{31.36}{3.61} & \bestmeanstd{73.02}{1.95} \\
CiteSeer / 1 & \meanstd{21.14}{2.87} & \meanstd{35.76}{6.80} & \meanstd{17.10}{2.64} & \meanstd{21.40}{4.60} & \meanstd{19.53}{3.34} & \bestmeanstd{47.12}{6.31} \\
CiteSeer / 3 & \meanstd{25.29}{2.16} & \meanstd{42.22}{5.88} & \meanstd{20.36}{2.02} & \meanstd{28.92}{2.68} & \meanstd{24.83}{2.77} & \bestmeanstd{61.20}{4.98} \\
CiteSeer / 5 & \meanstd{27.38}{2.84} & \meanstd{47.89}{4.85} & \meanstd{21.40}{2.57} & \meanstd{31.82}{2.60} & \meanstd{27.66}{2.49} & \bestmeanstd{66.17}{2.16} \\
PubMed / 1 & \meanstd{55.44}{11.42} & \meanstd{48.77}{8.54} & \meanstd{36.29}{5.79} & \meanstd{41.12}{6.83} & \meanstd{50.24}{12.95} & \bestmeanstd{58.70}{9.78} \\
PubMed / 3 & \meanstd{65.55}{8.11} & \meanstd{59.75}{7.97} & \meanstd{51.28}{6.93} & \meanstd{57.46}{4.96} & \meanstd{62.47}{5.54} & \bestmeanstd{68.66}{4.58} \\
PubMed / 5 & \meanstd{68.37}{3.41} & \meanstd{63.96}{5.18} & \meanstd{55.83}{4.93} & \meanstd{61.37}{5.32} & \meanstd{66.94}{4.09} & \bestmeanstd{70.98}{2.46} \\
\bottomrule
\end{tabular}
}
\end{table}

\begin{table}[!htbp]
\centering
\caption{\lang{Macro-F1 (\%, higher is better). Main comparison. Within each dataset/shot/method, training seeds are averaged within each support draw; mean and sample standard deviation are then computed over 20 support draws. Bold marks the largest mean in that row. $\dagger$: common-protocol adaptation.}{Macro-F1（\%，越高越好）。主比较。固定数据集、shot 与方法，在每个支持抽样内先平均训练 seed，再对 20 个支持抽样计算均值与样本标准差。粗体为该行最大均值。 $\dagger$：统一协议适配实现。}}
\label{tab:main-macro_f1}
\small
\setlength{\tabcolsep}{4pt}
\resizebox{\linewidth}{!}{%
\begin{tabular}{lrrrrrr}
\toprule
\lang{Dataset / shot}{数据集 / shot} & \frozenname & \ftname & GPF & GPF-plus & \gpromptname & \method \\
\midrule
Cora / 1 & \meanstd{22.51}{4.82} & \meanstd{44.38}{10.07} & \meanstd{7.08}{2.07} & \meanstd{27.08}{6.92} & \meanstd{21.03}{4.36} & \bestmeanstd{53.34}{8.94} \\
Cora / 3 & \meanstd{29.07}{3.47} & \meanstd{60.94}{6.41} & \meanstd{21.14}{6.64} & \meanstd{39.31}{4.55} & \meanstd{26.98}{3.41} & \bestmeanstd{69.38}{2.91} \\
Cora / 5 & \meanstd{33.07}{4.47} & \meanstd{66.99}{4.52} & \meanstd{25.87}{6.19} & \meanstd{51.82}{5.06} & \meanstd{30.82}{3.43} & \bestmeanstd{72.04}{1.77} \\
CiteSeer / 1 & \meanstd{20.39}{2.62} & \meanstd{30.99}{6.30} & \meanstd{6.78}{1.56} & \meanstd{12.83}{4.10} & \meanstd{18.00}{2.95} & \bestmeanstd{44.26}{5.75} \\
CiteSeer / 3 & \meanstd{24.32}{1.87} & \meanstd{38.43}{5.80} & \meanstd{13.81}{5.45} & \meanstd{26.16}{2.85} & \meanstd{23.71}{2.39} & \bestmeanstd{58.19}{4.67} \\
CiteSeer / 5 & \meanstd{26.45}{2.71} & \meanstd{44.62}{4.98} & \meanstd{15.27}{5.85} & \meanstd{29.46}{2.64} & \meanstd{26.72}{2.49} & \bestmeanstd{63.03}{2.12} \\
PubMed / 1 & \meanstd{53.44}{10.49} & \meanstd{44.52}{9.85} & \meanstd{25.62}{6.56} & \meanstd{34.89}{8.06} & \meanstd{45.35}{10.95} & \bestmeanstd{57.37}{10.07} \\
PubMed / 3 & \meanstd{63.19}{8.60} & \meanstd{57.64}{9.26} & \meanstd{46.88}{9.09} & \meanstd{54.47}{4.51} & \meanstd{59.93}{5.76} & \bestmeanstd{68.21}{4.08} \\
PubMed / 5 & \meanstd{66.27}{3.55} & \meanstd{62.53}{6.05} & \meanstd{52.42}{6.87} & \meanstd{58.79}{5.19} & \meanstd{64.79}{4.02} & \bestmeanstd{70.39}{2.47} \\
\bottomrule
\end{tabular}
}
\end{table}


Cora 的 5-shot Accuracy 为 73.02\%，冻结预训练与全量微调分别为 33.56\% 和 67.50\%；CiteSeer 在相同 shot 数下分别为 66.17\%、27.38\%、47.89\%。PubMed 的冻结预训练更强：5-shot 为 68.37\%，\method 为 70.98\%，全量微调为 63.96\%。因此，相对冻结预训练的改善在 PubMed 上明显较小，跨数据集平均会掩盖这一依赖关系。

逐条件比较中，\method 的 Accuracy 与 Macro-F1 也均高于三类提示算子。针对五个主比较方法，两项指标的配对区间全部为正。这些区间刻画当前协议下的 seed 波动，不能推出对所有图提示方法或独立图总体的普遍优势。完整配对差值见附录~\ref{app:paired}。

\subsection{原型路由是否优于共享提示？}
核心消融将式~\eqref{eq:injection} 的 $B$ 替换为元素均为 $1/C$ 的常数矩阵，保留提示矩阵、尺度、注入位置、优化与最终读出，移除节点相关的提示分配。常数矩阵直接用于组合，不再中心化，否则提示会被完全消除；同时保留无提示的冻结预训练条件。
\begin{table}[!htbp]
\centering
\caption{\lang{Accuracy (\%, higher is better). Prototype-routing ablation. Within each dataset/shot/method, training seeds are averaged within each support draw; mean and sample standard deviation are then computed over 20 support draws. Bold marks the largest mean in that row.}{Accuracy（\%，越高越好）。原型路由消融。固定数据集、shot 与方法，在每个支持抽样内先平均训练 seed，再对 20 个支持抽样计算均值与样本标准差。粗体为该行最大均值。}}
\label{tab:ablation-accuracy}
\small
\setlength{\tabcolsep}{4pt}
\begin{tabular}{lrrr}
\toprule
\lang{Dataset / shot}{数据集 / shot} & \frozenname & \sharedname & \method \\
\midrule
Cora / 1 & \meanstd{23.17}{5.48} & \meanstd{22.62}{4.41} & \bestmeanstd{54.59}{9.18} \\
Cora / 3 & \meanstd{29.42}{3.47} & \meanstd{31.70}{4.95} & \bestmeanstd{70.76}{2.76} \\
Cora / 5 & \meanstd{33.56}{5.03} & \meanstd{36.08}{4.19} & \bestmeanstd{73.02}{1.95} \\
CiteSeer / 1 & \meanstd{21.14}{2.87} & \meanstd{20.72}{2.46} & \bestmeanstd{47.12}{6.31} \\
CiteSeer / 3 & \meanstd{25.29}{2.16} & \meanstd{24.01}{2.70} & \bestmeanstd{61.20}{4.98} \\
CiteSeer / 5 & \meanstd{27.38}{2.84} & \meanstd{25.73}{2.78} & \bestmeanstd{66.17}{2.16} \\
PubMed / 1 & \meanstd{55.44}{11.42} & \meanstd{53.84}{10.55} & \bestmeanstd{58.70}{9.78} \\
PubMed / 3 & \meanstd{65.55}{8.11} & \meanstd{62.96}{7.94} & \bestmeanstd{68.66}{4.58} \\
PubMed / 5 & \meanstd{68.37}{3.41} & \meanstd{67.95}{3.11} & \bestmeanstd{70.98}{2.46} \\
\bottomrule
\end{tabular}
\end{table}

\begin{table}[htbp]
\centering
\caption{\lang{Macro-F1 (\%, higher is better). Prototype-routing ablation. Within each dataset/shot/method, training seeds are averaged within each support draw; mean and sample standard deviation are then computed over 20 support draws. Bold marks the largest mean in that row.}{Macro-F1（\%，越高越好）。原型路由消融。固定数据集、shot 与方法，在每个支持抽样内先平均训练 seed，再对 20 个支持抽样计算均值与样本标准差。粗体为该行最大均值。}}
\label{tab:ablation-macro_f1}
\small
\setlength{\tabcolsep}{4pt}
\resizebox{\linewidth}{!}{%
\begin{tabular}{lrrr}
\toprule
\lang{Dataset / shot}{数据集 / shot} & \frozenname & \sharedname & \method \\
\midrule
Cora / 1 & \meanstd{22.51}{4.82} & \meanstd{21.65}{3.66} & \bestmeanstd{53.34}{8.94} \\
Cora / 3 & \meanstd{29.07}{3.47} & \meanstd{30.25}{4.85} & \bestmeanstd{69.38}{2.91} \\
Cora / 5 & \meanstd{33.07}{4.47} & \meanstd{34.89}{4.10} & \bestmeanstd{72.04}{1.77} \\
CiteSeer / 1 & \meanstd{20.39}{2.62} & \meanstd{19.57}{2.33} & \bestmeanstd{44.26}{5.75} \\
CiteSeer / 3 & \meanstd{24.32}{1.87} & \meanstd{22.67}{2.46} & \bestmeanstd{58.19}{4.67} \\
CiteSeer / 5 & \meanstd{26.45}{2.71} & \meanstd{24.55}{2.74} & \bestmeanstd{63.03}{2.12} \\
PubMed / 1 & \meanstd{53.44}{10.49} & \meanstd{51.93}{9.81} & \bestmeanstd{57.37}{10.07} \\
PubMed / 3 & \meanstd{63.19}{8.60} & \meanstd{60.94}{8.60} & \bestmeanstd{68.21}{4.08} \\
PubMed / 5 & \meanstd{66.27}{3.55} & \meanstd{66.29}{2.87} & \bestmeanstd{70.39}{2.47} \\
\bottomrule
\end{tabular}
}
\end{table}


表~\ref{tab:ablation-accuracy} 与表~\ref{tab:ablation-macro_f1} 表明，单纯学习共享提示不足以解释完整方法的收益。5-shot 下，原型路由相对共享提示在 Cora、CiteSeer、PubMed 的 Accuracy 分别提高 36.94、40.45、3.02 个百分点，对应配对 95\% 区间为 $[34.87,39.06]$、$[39.06,41.84]$、$[1.39,4.51]$ 个百分点。在各个 1-shot 和 3-shot 条件中，两项主要指标的配对改善也均为正。

该比较支持完整原型路由构造相对共享提示的作用，不能将收益单独归因于原型校正、中心化或图聚合中的某一步。均匀组合虽然保留名义提示参数量，却降低有效自由度，因此也不能视为完全控制容量后对唯一路由机制的证明。

\section{局限性}
实验覆盖三个引文网络，每个图仅使用一份固定预训练权重，任务为同图已知类别分类。预训练 seed 未知，未估计不同预训练运行之间的波动。报告范围在实验完成后确定，因此属于描述性分析，不是新的独立确认实验；结果不支持对图分类、未见类别或任意预训练来源的推广结论。

统一协议下的提示基线及有限选参范围未必达到各方法最佳已发表成绩。基线选参产生的训练预算也不同于 \method 的固定配方。路由消融改变了有效提示容量，未区分路由内部各步骤的贡献。此外，方法利用查询特征和全图统计，缺少这类传导式信息的部署条件需另行评估。

\section{结论}
\method 利用支持集锚定的原型组织可学习隐藏层提示的分配，从而适配冻结图编码器。在所评估的 Cora、CiteSeer、PubMed 条件下，逐条件主比较与共享提示消融共同支持这种分配方式的价值。现有证据支持原型路由作为当前协议下的适配设计，其对其他图类型和预训练来源的普适性仍需进一步验证。

\section*{潜在影响}
本文研究公开引文网络基准，不评估涉及个人的实际决策。图聚合与未标注原型校正可能传播错误，尤其当图结构或类别结构偏离所评估设置时。涉及重要后果的应用需要开展基准之外的领域评估。
